% Interpolation theorem for weighted Chebyshev races — DRAFT FOR RED TEAM
% To be merged into weighted-races.tex as a new section once verified.
% Compile standalone with the weighted-races.tex preamble if needed.
%
% STATUS: draft; every step marked [ANS], [RS], [Hum], [Dev] leans on the
% cited work — the red team must check each borrowed hypothesis against
% the actual papers.

\section{An interpolation theorem}\label{sec:theorem}

Throughout this section $\chi=\chi_4$, $\rho=\tfrac12+i\gamma$ runs over
the nontrivial zeros of $L(s,\chi)$ with $\gamma>0$, GRH$(\chi)$ denotes
the Riemann hypothesis for $L(s,\chi)$, and LI$(\chi)$ the linear
independence over $\mathbb Q$ of the positive ordinates $\gamma$. For
$m\in(-\tfrac12,0]$ define the ($\theta$-form) weighted race
\[
R_m(x)\;=\;\sum_{\substack{p\le x\\ p\equiv3\,(4)}}p^m\log p
\;-\;\sum_{\substack{p\le x\\ p\equiv1\,(4)}}p^m\log p
\;=\;-\sum_{p\le x}\chi(p)\,p^m\log p ,
\]
and its normalization
\[
E_m(x)\;=\;(2m+1)\,x^{-(m+1/2)}\,R_m(x).
\]
Write $a_\gamma(m)=(2m+1)\bigl((m+\tfrac12)^2+\gamma^2\bigr)^{-1/2}$ and
\[
X_m\;=\;1+\sum_{\gamma>0}2a_\gamma(m)\cos\theta_\gamma ,
\qquad \theta_\gamma\ \text{i.i.d.\ uniform on }[0,2\pi),
\]
the series converging a.s.\ and in $L^2$ since
$\sum_\gamma a_\gamma(m)^2<\infty$ (zero-counting density
$\ll\log\gamma$). Set $C:=\sum_{\gamma>0}2\gamma^{-2}$
($=0.156033\ldots$ numerically).

\begin{theorem}\label{thm:interp}
Assume GRH$(\chi_4)$.
\begin{enumerate}
\item[(i)] For each $m\in(-\tfrac12,0]$ the function
$u\mapsto E_m(e^u)$ possesses a limiting logarithmic distribution
$\nu_m$: for every bounded Lipschitz $f$,
$\lim_{Y\to\infty}\frac1Y\int_0^Y f(E_m(e^u))\,du
=\int_{\mathbb R}f\,d\nu_m$. Under LI$(\chi_4)$, $\nu_m$ is the law of
$X_m$.
\item[(ii)] Under LI$(\chi_4)$ the density
$\delta(m):=\nu_m\bigl((0,\infty)\bigr)$ satisfies
$\tfrac12<\delta(m)<1$ for every $m\in(-\tfrac12,0]$, and
$m\mapsto\delta(m)$ is continuous on $(-\tfrac12,0]$.
\item[(iii)] Under LI$(\chi_4)$, for all $m\in(-\tfrac12,0]$,
\[
1-\delta(m)\;\le\;C\,(2m+1)^2 ,
\qquad\text{and in fact}\qquad
1-\delta(m)\;\le\;\exp\!\Bigl(-\frac{1}{8C\,(2m+1)^2}\Bigr)
\ \text{for }(2m+1)^2\le(8C)^{-1}.
\]
In particular $\delta(m)\to1$ as $m\to-\tfrac12^+$.
\item[(iv)] Consequently the classical density formulation of
Chebyshev's bias ($m=0$: $\delta(0)=0.995928\ldots$
\emph{[RS]}) and the weighted formulation at the threshold
($m=-\tfrac12$: $R_{-1/2}$ of one sign for large $x$ with
$\tfrac12\log\log x$ drift \emph{[Aoki--Koyama; Sheth]}) are endpoints
of the single continuous family $\delta(m)$, with maximal bias
$\delta\to1$ attained exactly in the limit toward the weighted
threshold.
\end{enumerate}
\end{theorem}

\begin{remark}
Part (iii) is the arithmetic-progression analogue of Humphries's
$\lim_{\alpha\to1/2^-}\delta(P_\alpha)=1$ for weighted sums of the
Liouville function \emph{[Hum, Thm.~1.5]}; the proof strategy
(variance collapse against a fixed bias) is the same, and we claim
novelty only for the statement in the prime-race setting together with
the explicit constant. Part (i) does not require LI; this follows the
framework of Akbary--Ng--Shahabi \emph{[ANS]} (cf.\ Devin \emph{[Dev]}).
\end{remark}

\begin{proof}[Proof of (i)]
Under GRH$(\chi)$ the explicit formula for
$\psi_m(x):=\sum_{n\le x}\Lambda(n)\chi(n)n^m$ gives, for
$2\le T\le x$,
\begin{equation}\label{eq:trunc}
\psi_m(x)\;=\;-\sum_{|\gamma|\le T}\frac{x^{\rho+m}}{\rho+m}
\;+\;O\!\Bigl(\frac{x^{m+1}\log^2(xT)}{T}+x^{m}\log x\Bigr),
\end{equation}
obtained from the classical truncated formula for $\psi(x,\chi)$
by partial summation ($\psi_m(x)=x^m\psi(x,\chi)-m\int_2^x
t^{m-1}\psi(t,\chi)\,dt$) and the corresponding truncated expansion of
$\psi(t,\chi)$ inserted in the integral; the shift replaces each
$x^\rho/\rho$ by $x^{\rho+m}/(\rho+m)$ since
$x^m\frac{x^\rho}{\rho}-m\int_2^x t^{m-1}\frac{t^\rho}{\rho}dt
=\frac{x^{\rho+m}}{\rho+m}+O_\gamma(1)$.
Separating prime powers,
\[
-R_m(x)=\psi_m(x)-\sum_{p\le\sqrt x}p^{2m}\log p+O\!\bigl(x^{(2m+1)/3}\log x\bigr),
\]
and $\sum_{p\le\sqrt x}p^{2m}\log p=\frac{x^{m+1/2}}{2m+1}
+O\bigl(x^{m+1/2}e^{-c\sqrt{\log x}}\bigr)$ by the prime number theorem
with classical error term and partial summation. Note
$(2m+1)/3<m+\tfrac12$ exactly when $m>-\tfrac12$, so both error terms
are $o(x^{m+1/2})$ for fixed $m\in(-\tfrac12,0]$. Hence with $x=e^u$,
\[
E_m(e^u)\;=\;1+\mathrm{Re}\sum_{0<\gamma\le T} r_\gamma(m)\,e^{i\gamma u}
+\;\mathcal E(u,T),\qquad
r_\gamma(m)=\frac{2(2m+1)}{m+\tfrac12+i\gamma},
\]
where, by \eqref{eq:trunc} with $T=e^{u/2}$ (say) and the mean-square
bound of \emph{[ANS, \S2--4]} for the tail of the zero sum,
$\limsup_{Y\to\infty}\frac1Y\int_0^Y|\mathcal E(u,T)|^2du\to0$ as
$T\to\infty$. Since $\sum_\gamma|r_\gamma(m)|^2<\infty$, the hypotheses
of the limiting-distribution theorem of \emph{[ANS, Thm.~1.2]} hold for
the shifted amplitudes exactly as in their unshifted applications
($\mathrm{Re}(\rho+m)=m+\tfrac12$ is constant along the zeros, playing
the role their $\tfrac12$ plays), and (i) follows; the identification
of $\nu_m$ as the law of $X_m$ under LI is \emph{[ANS, \S3]} (or
\emph{[RS, \S2]}) verbatim.
\end{proof}

\begin{proof}[Proof of (ii)]
Fix $m\in(-\tfrac12,0]$. The random variable $X_m-1$ is a sum of
independent symmetric variables $2a_\gamma\cos\theta_\gamma$ with
arcsine (hence absolutely continuous) laws; the sum of an absolutely
continuous independent component with anything is absolutely
continuous, so $X_m$ has a density and in particular no atom at $0$.
Since $\sum_\gamma a_\gamma(m)=\infty$ (by the zero-counting theorem
$N(T,\chi)\asymp T\log T$, $\sum_\gamma\gamma^{-1}$ diverges), the
support of $X_m$ is all of $\mathbb R$ \emph{[RS, proof of positivity
of all densities]}; combined with symmetry of $X_m-1$ about $0$ and
$\mathbb P(X_m>0)\ge\mathbb P(X_m-1>0)=\tfrac12$ plus full support,
$\tfrac12<\delta(m)<1$.

For continuity: the characteristic function of $X_m-1$ is
$\Phi_m(t)=\prod_{\gamma}J_0\bigl(2a_\gamma(m)t\bigr)$. On a compact
$K\subset(-\tfrac12,0]$ and compact $t$-sets, $\log J_0(2a_\gamma t)
=-a_\gamma^2t^2+O(a_\gamma^4t^4)$ uniformly for $\gamma$ large, and
$\sum_\gamma a_\gamma(m)^2$ is continuous in $m$ on $K$ (dominated
convergence against $\sum(2m+1)^2/\gamma^2$); hence
$\Phi_{m'}\to\Phi_m$ pointwise as $m'\to m$, so
$X_{m'}\Rightarrow X_m$ weakly (L\'evy). Since $X_m$ has no atom at
$0$, $\delta(m')=\mathbb P(X_{m'}>0)\to\mathbb P(X_m>0)=\delta(m)$.
\end{proof}

\begin{proof}[Proof of (iii)]
Let $V_m=X_m-1$, $\mathbb E V_m=0$,
$\mathrm{Var}(V_m)=\sum_\gamma 2a_\gamma(m)^2
\le(2m+1)^2\sum_\gamma\frac{2}{\gamma^2}=C(2m+1)^2$.
By Chebyshev, $1-\delta(m)\le\mathbb P(V_m\le-1)\le C(2m+1)^2$.
For the exponential form: the summands $2a_\gamma\cos\theta_\gamma$
are independent, mean zero, and bounded by $b_\gamma=2a_\gamma$ with
$\sum_\gamma b_\gamma^2=4\sum a_\gamma^2\le 2C(2m+1)^2\cdot 2
=4C(2m+1)^2$; wait --- $\sum b_\gamma^2 = 4\sum a_\gamma^2
\le 4\cdot\tfrac{C}{2}(2m+1)^2\cdot ?$. We record it cleanly:
$\sum_\gamma a_\gamma^2\le\tfrac{C}{2}(2m+1)^2$ (since
$\sum 2a_\gamma^2\le C(2m+1)^2$), so $\sum b_\gamma^2\le2C(2m+1)^2$.
Hoeffding's inequality for sums of bounded independent variables gives
\[
\mathbb P(V_m\le-1)\;\le\;\exp\!\Bigl(-\frac{2\cdot 1^2}{4\sum_\gamma
a_\gamma^2\cdot ?}\Bigr)
\]
--- \textbf{[RED TEAM: fix the Hoeffding constant.]} With the standard
form $\mathbb P(S\le-t)\le\exp(-t^2/(2\sum c_\gamma^2))$ for summands
in $[-c_\gamma,c_\gamma]$... the summand $2a_\gamma\cos\theta$ lies in
$[-2a_\gamma,2a_\gamma]$, i.e.\ range $4a_\gamma$; Hoeffding
(range form) gives $\mathbb P(S\le -t)\le\exp\bigl(-2t^2/\sum(\text{range})^2\bigr)
=\exp\bigl(-2/(16\sum a_\gamma^2)\bigr)
=\exp\bigl(-1/(8\sum a_\gamma^2)\bigr)
\ge\exp\bigl(-1/(4C(2m+1)^2)\bigr)$ --- so
$1-\delta(m)\le\exp\bigl(-1/(4C(2m+1)^2)\bigr)$,
valid for all $m$, which is stronger than stated in the theorem;
\textbf{the red team should fix the final constant and adjust the
theorem statement to whatever is correct.}
Both bounds give $\delta(m)\to1$ as $m\to-\tfrac12^+$.
\end{proof}

\begin{proof}[Proof of (iv)]
Immediate from (i)--(iii), \emph{[RS]} for $\delta(0)$, and
\emph{[Aoki--Koyama; Sheth]} for the $m=-\tfrac12$ endpoint (whose
statement requires no LI; under DRH resp.\ GRH off a set of finite
logarithmic measure).
\end{proof}

\begin{remark}[The unlogged race]
The same statements hold for the unlogged races
$D_m(x)=\sum_{p\equiv3}p^m-\sum_{p\equiv1}p^m$ normalized by
$(2m+1)(m+\tfrac12)^{-1}\cdots$ --- more precisely, partial summation
$D_m(x)=R_m(x)/\log x+\int_2^x R_m(t)\,dt/(t\log^2t)$ multiplies every
term of the expansion (bias and each zero term) by the same factor
$(\log x)^{-1}(1+O(1/\log x))$ up to $o(x^{m+1/2}/\log x)$, so the
normalized limiting distribution is unchanged. \textbf{[RED TEAM:
verify this lemma's error bookkeeping --- it was checked numerically
(corrections 0.70\% vs 0.70\% predicted at $x=10^{10}$) but the
$o(\cdot)$ claim needs the same mean-square treatment as (i).]}
\end{remark}

% Bibliography keys to merge: [ANS]=Akbary-Ng-Shahabi QJM 65 (2014);
% [RS]=Rubinstein-Sarnak 1994; [Hum]=Humphries JNT 133 (2013);
% [Dev]=Devin MPCPS 169 (2020); [Aoki-Koyama]=JNT 245 (2023);
% [Sheth]=MPCPS 179 (2025).
